\documentclass[11pt]{article}

\usepackage[margin=1in]{geometry}
\usepackage{amsmath, amssymb}
\usepackage{xcolor}
\usepackage{listings}
\usepackage{hyperref}
\usepackage{array}
\usepackage{booktabs}

\hypersetup{
    colorlinks=true,
    linkcolor=blue!60!black,
    urlcolor=blue!60!black,
    citecolor=blue!60!black
}

\lstdefinestyle{bugpython}{
  language=Python,
  basicstyle=\ttfamily\small,
  keywordstyle=\color{blue!60!black},
  stringstyle=\color{green!40!black},
  commentstyle=\color{gray!70!black},
  showstringspaces=false,
  breaklines=true,
  frame=single,
  rulecolor=\color{gray!30},
  columns=fullflexible,
  keepspaces=true
}

\title{SymPy Bug Report}
\author{}
\date{}

\begin{document}

\maketitle

\section{Bug 32: solveset misses a closed endpoint for \texorpdfstring{\(\exp(i x)=1\)}{exp(i*x)=1}}

\subsection{Minimal Reproducer}

\medskip

\begin{lstlisting}[style=bugpython]
from sympy import Eq, I, Interval, exp, pi, simplify, solveset, symbols

x = symbols("x")
interval = Interval(0, 2*pi)
sol = solveset(Eq(exp(I*x), 1), x, interval)

print("solution:", sol)
print("2*pi in solution:", sol.contains(2*pi))
print("residual at 2*pi:", simplify(exp(I*2*pi) - 1))
\end{lstlisting}

\subsection{Observed Behavior}

\medskip

\begin{lstlisting}[style=bugpython]
solution: {0}
2*pi in solution: False
residual at 2*pi: 0
\end{lstlisting}

SymPy returns only \(\{0\}\).  The endpoint \(2\pi\) is reported as absent from
the solution set even though direct simplification of the equation residual at
that point gives \(0\).

\subsection{Expected Behavior}

On the closed interval \([0,2\pi]\), the equation
\[
    \exp(i x)=1
\]
has two endpoint solutions, \(x=0\) and \(x=2\pi\).  The expected result is
\[
    \{0,2\pi\}.
\]

\subsection{Mathematical Explanation}

For real \(x\), Euler's formula gives
\[
    \exp(i x)=\cos x+i\sin x.
\]
This equals \(1\) exactly when \(\cos x=1\) and \(\sin x=0\), i.e. at
\[
    x=2\pi k,\qquad k\in\mathbb{Z}.
\]
Intersecting this periodic solution family with the closed interval
\([0,2\pi]\) gives
\[
    \{2\pi k:k\in\mathbb{Z}\}\cap[0,2\pi]=\{0,2\pi\}.
\]
The missing endpoint is therefore not a harmless choice of representative or an
equivalent periodic notation: \(\{0\}\) is a strict subset of the correct
closed-interval solution set.

\subsection{Source-Level Diagnosis}

\begin{sloppypar}
The diagnosis identifies the relevant implementation as
\nolinkurl{sympy/solvers/solveset.py}.  The public call
\texttt{solveset(Eq(exp(I*x), 1), x, Interval(0, 2*pi))} is rewritten as the
zero equation \(\exp(i x)-1=0\).  Since the requested domain is a subset of the
reals, \texttt{solveset} creates a real dummy symbol and reaches the generic
inversion branch in \texttt{\_solveset}, where it calls
\texttt{inverter(f, 0, symbol)}.  That inverter is \texttt{\_invert}, and the
diagnosis locates the problematic dispatch at
\nolinkurl{sympy/solvers/solveset.py:178}.
\end{sloppypar}

\begin{lstlisting}[style=bugpython]
if domain.is_subset(S.Reals):
    x1, s = _invert_real(f_x, FiniteSet(y), x)
else:
    x1, s = _invert_complex(f_x, FiniteSet(y), x)
\end{lstlisting}

This dispatch chooses \texttt{\_invert\_real} based on the solution domain
rather than on the structure of the expression being inverted.  For
\(\exp(i x)\), the expression has complex periodic behavior even when \(x\) is
real.  The real exponential inversion path then applies a single-valued
logarithm rule for \texttt{exp}, reducing \(\exp(i x)=1\) to \(i x=0\) and
therefore \(x=0\).  By contrast, the diagnosis records that
\texttt{invert\_complex(exp(I*x), 1, x)} gives the periodic family
\texttt{ImageSet(Lambda(\_n, 2*\_n*pi), Integers)}.

The final domain intersection cannot recover the lost integer parameter once
the inversion step has collapsed the family to \(\{0\}\).  A plausible fix
direction identified by the diagnosis is for equations involving \texttt{exp}
with non-real or periodic arguments, especially \(\exp(i x)\) with real \(x\),
to use complex exponential inversion and then intersect the resulting periodic
image set with the real domain.  More generally, the \texttt{\_invert}
dispatch appears to need to consider the expression's value and argument
structure, not only \texttt{domain.is\_subset(S.Reals)}.

\subsection{Additional Failing Instances}

The independent verification checks the shifted family
\[
    \operatorname{solveset}\bigl(\exp(i(x+a))=1,\ x,\ [-a,2\pi-a]\bigr),
    \qquad a\in\{0,1,2,3,4,5\}.
\]
For each \(a\), the true solution family is \(x=2\pi k-a\).  The chosen closed
interval contains both endpoint solutions \(-a\) and \(2\pi-a\), but SymPy
returns only the left endpoint.

\begin{center}
\small
\renewcommand{\arraystretch}{1.3}
\begin{tabular}{@{}>{\raggedright\arraybackslash}p{0.44\linewidth}>{\raggedright\arraybackslash}p{0.23\linewidth}>{\raggedright\arraybackslash}p{0.23\linewidth}@{}}
\toprule
\textbf{Input / Expression} & \textbf{SymPy behavior} & \textbf{Expected behavior} \\
\midrule
\(\exp(i x)=1\) on \([0,2\pi]\) & returns \(\{0\}\); \(2\pi\notin\) result & include \(\{0,2\pi\}\) \\
\(\exp(i(x+1))=1\) on \([-1,2\pi-1]\) & returns \(\{-1\}\); \(2\pi-1\notin\) result & include \(\{-1,2\pi-1\}\) \\
\(\exp(i(x+2))=1\) on \([-2,2\pi-2]\) & returns \(\{-2\}\); \(2\pi-2\notin\) result & include \(\{-2,2\pi-2\}\) \\
\(\exp(i(x+3))=1\) on \([-3,2\pi-3]\) & returns \(\{-3\}\); \(2\pi-3\notin\) result & include \(\{-3,2\pi-3\}\) \\
\(\exp(i(x+4))=1\) on \([-4,2\pi-4]\) & returns \(\{-4\}\); \(2\pi-4\notin\) result & include \(\{-4,2\pi-4\}\) \\
\(\exp(i(x+5))=1\) on \([-5,2\pi-5]\) & returns \(\{-5\}\); \(2\pi-5\notin\) result & include \(\{-5,2\pi-5\}\) \\
\bottomrule
\end{tabular}
\end{center}

The expected behavior is the same across the family because
\(\exp(i(x+a))=1\) is equivalent, over real \(x\), to
\[
    x+a=2\pi k,\qquad k\in\mathbb{Z}.
\]
Thus both endpoints of each listed interval are exact solutions.

\subsection{Related Issues Assessment}

The best-effort deduplication check found no public match for this exact
\texttt{solveset} call, the exact returned set \(\{0\}\), or the missing
closed-interval endpoint \(2\pi\).  The closest conceptual relationship is the
broad family of periodic equation solving over restricted domains, but the
available evidence does not identify a public issue, PR, release-note entry,
Stack Overflow post, or mailing-list discussion that clearly covers this case.
The deduplication verdict is therefore a new failure mode with no public
precedent, and the bug remains individually actionable as a concise
\texttt{solveset} correctness regression test.

\subsection{Proposed SymPy Regression Test}

\medskip

\begin{lstlisting}[style=bugpython]
import pytest
from sympy import Eq, I, Interval, S, exp, pi, solveset, symbols


@pytest.mark.parametrize("a", [0, 1, 2, 3, 4, 5])
def test_solveset_exp_i_x_includes_closed_period_endpoint(a):
    x = symbols("x")
    sol = solveset(Eq(exp(I*(x + a)), 1), x, Interval(-a, 2*pi - a))
    assert sol.contains(2*pi - a) is S.true
\end{lstlisting}

\end{document}
